\section{Po co teoria kategorii?}\label{sec:motywacja}

W wielu działach matematyki pojawia się podobna sytuacja. Mamy pewnego rodzaju obiekty i odwzorowania, które zachowują strukturę tych obiektów. Dla zbiorów są to funkcje, dla grup homomorfizmy, dla przestrzeni topologicznych odwzorowania ciągłe, a dla przestrzeni wektorowych odwzorowania liniowe. Odwzorowania można składać, a każdy obiekt ma odwzorowanie identycznościowe.

Teoria kategorii bada własności, które można sformułować wyłącznie za pomocą obiektów, odwzorowań i ich składania. Nie usiłuje zastąpić algebry lub topologii. Pozwala natomiast rozpoznawać te same konstrukcje w różnych działach matematyki i wyjaśnia, dlaczego są one do siebie podobne.

\begin{przyklad}[Iloczyn kartezjański i iloczyn grup]
Dla zbiorów $A,B$ definiujemy $A\times B$ i rzuty $p_A:A\times B\to A$, $p_B:A\times B\to B$. Jeśli mamy funkcje $f:X\to A$ i $g:X\to B$, istnieje dokładnie jedna funkcja $h:X\to A\times B$ spełniająca $p_A\circ h=f$ oraz $p_B\circ h=g$, mianowicie $h(x)=(f(x),g(x))$. Ta sama własność charakteryzuje iloczyn grup z rzutami będącymi homomorfizmami. Rzeczywiście, dla homomorfizmów $f:X\to A$ i $g:X\to B$ oraz $x,y\in X$ mamy
\[
h(xy)=(f(xy),g(xy))=(f(x)f(y),g(x)g(y))=h(x)h(y).
\]
Zatem $h$ jest homomorfizmem. Jeżeli inny homomorfizm $u:X\to A\times B$ ma te same rzuty, to dla każdego $x\in X$ obie współrzędne $u(x)$ są równe odpowiednim współrzędnym $h(x)$, więc $u=h$. Własność tę nazwiemy \emph{uniwersalną}. Jej sformułowanie, w przeciwieństwie do dowodu istnienia, nie odwołuje się do elementów iloczynu. Rozwiniemy ją w rozdziale~\ref{sec:produkty}.
\end{przyklad}

\begin{przyklad}[Ta sama własność w zbiorze uporządkowanym]\label{cat:ex:motywacja-porzadek}
Rozważmy zbiór częściowo uporządkowany $(P,\leq)$. Między $a$ i $b$ dopuszczamy dokładnie jeden morfizm $a\to b$, jeśli $a\leq b$, a w przeciwnym razie żadnego. Zwrotność zapewnia identyczności, przechodniość pozwala składać morfizmy, a łączność wynika z tego, że między ustalonymi obiektami istnieje co najwyżej jeden morfizm.

Odpowiednikiem iloczynu $a$ i $b$ jest ich kres dolny $m$, o ile istnieje. Istnienie rzutów $m\to a$ i $m\to b$ oznacza $m\leq a$ i $m\leq b$. Jeżeli $x\leq a$ oraz $x\leq b$, własność uniwersalna wymaga dokładnie jednego morfizmu $x\to m$, czyli $x\leq m$. Są to dokładnie warunki definiujące największe ograniczenie dolne. W $\mathbb R$ z naturalnym porządkiem iloczynem w tym sensie jest $\min(a,b)$, a nie zwykłe mnożenie liczb.
\end{przyklad}

\begin{przyklad}[Dlaczego wybór morfizmów ma znaczenie?]\label{cat:ex:motywacja-morfizmy}
Rozważmy wszystkie zbiory. Jeśli jako morfizmy dopuszczamy wszystkie funkcje, każdy singleton $S$ jest obiektem, do którego z dowolnego zbioru $X$ prowadzi dokładnie jeden morfizm: funkcja przyjmująca jedyną możliwą wartość (dla $X=\varnothing$ jest to funkcja pusta).

Jeżeli dopuścimy tylko funkcje różnowartościowe, identyczności i złożenia nadal są dozwolone, lecz ze zbioru dwuelementowego do singletonu nie istnieje żaden morfizm. Co więcej, w tej kategorii nie istnieje obiekt, do którego z każdego zbioru prowadzi dokładnie jeden morfizm. Taki obiekt musiałby mieć dokładnie jeden element, ponieważ morfizmy z singletonu odpowiadają jego elementom. Tymczasem ze zbioru dwuelementowego nie ma do niego iniekcji. Pojęcie obiektu końcowego omówimy w rozdziale~\ref{sec:uniwersalne}.
\end{przyklad}

\begin{uwaga}[Jak czytać diagramy]
Zapis $A\xrightarrow{f}B$ oznacza morfizm $f$ z $A$ do $B$. Diagram jest \emph{przemienny}, gdy dowolne dwie skierowane drogi prowadzące od tego samego obiektu do tego samego obiektu końcowego wyznaczają ten sam morfizm. Przykładowo diagram
\begin{center}
\begin{tikzpicture}[>=stealth]
 \node (a) at (0,2) {$A$};
 \node (b) at (4,2) {$B$};
 \node (c) at (0,0) {$C$};
 \node (d) at (4,0) {$D$};
 \draw[->] (a) -- node[above] {$f$} (b);
 \draw[->] (a) -- node[left] {$g$} (c);
 \draw[->] (b) -- node[right] {$h$} (d);
 \draw[->] (c) -- node[below] {$k$} (d);
\end{tikzpicture}
\end{center}
jest przemienny wtedy i tylko wtedy, gdy $h\circ f=k\circ g$. Diagram jest więc graficznym zapisem równania. W $h\circ f$ najpierw wykonujemy $f$, a potem $h$. Sam rysunek kwadratu nie zapewnia przemienności; trzeba ją założyć lub wykazać.
\end{uwaga}

